% Generated draft: requires mathematical and source review.
\documentclass[11pt]{article}
\usepackage[margin=25mm]{geometry}
\usepackage{fontspec}
\setmainfont{Latin Modern Roman}
\usepackage{amsmath,amssymb,mathtools,unicode-math}
\setmathfont{Latin Modern Math}
\usepackage{enumitem,needspace,xurl,hyperref}
\setlength{\parindent}{0pt}
\setlength{\parskip}{5pt}
\setlength{\emergencystretch}{3em}
\newcommand{\N}{\mathbb{N}}
\newcommand{\Z}{\mathbb{Z}}
\newcommand{\Q}{\mathbb{Q}}
\newcommand{\R}{\mathbb{R}}
\newcommand{\C}{\mathbb{C}}
\newcommand{\F}{\mathbb{F}}
\newcommand{\A}{\mathbb{A}}
\newcommand{\PP}{\mathbb{P}}
\newcommand{\E}{\mathbb{E}}
\newcommand{\eps}{\varepsilon}
\newcommand{\dd}{\,\mathrm{d}}
\begin{document}
\section*{UnsolvedMath ID: 6900004}
\textbf{Source problem code:} catalogue entry
\subsection{Definitions and mathematical statement}
Let \(K_n\) be the complete graph on labeled vertices \(\{1,\ldots,n\}\). For a graph \(G\subseteq K_n\) and a placement \(p=(p_1,\ldots,p_n)\) in \(\R^d\), its equilibrium self-stress space is
\[
W(G,p)=\left\{(\omega_{ij})_{\{i,j\}\in E(G)}:
\forall i,\quad\sum_{j:\{i,j\}\in E(G)}\omega_{ij}(p_j-p_i)=0\right\},
\]
with one real coefficient per undirected edge. This supplies the tensegrity-fiber framework. Stratification can record more than the dimension of \(W(G,p)\), for example its admissible sign patterns and connectedness of corresponding placement sets; those conventions are not given in the extraction.

For a fixed common admissible placement space \(B_{d,n}\) and the source's stratification \(\mathscr S_d(G)\), the requested graph-equivalence classification is to characterize
\[
G\equiv_d G'\quad\Longleftrightarrow\quad
\mathscr S_d(G)=\mathscr S_d(G').
\]
Literal equality of partitions means
\[
\forall p,q\in B_{d,n},\qquad
p,q\text{ lie in one }G\text{-stratum}
\iff p,q\text{ lie in one }G'\text{-stratum}.
\]
This common-base formulation regards each graph as an edge set on the same \(n\) labels, adding isolated vertices if that is the intended convention. It does not identify graph isomorphism with equality of labeled placement strata or impose only a rank-stratification interpretation.

The question asks which pairs or equivalence classes of subgraphs satisfy this condition. The ambient dimension, allowed placements, quotient by Euclidean/affine motions, precise stratum equivalence, and literal equality versus equivalence under a base homeomorphism are not specified. The supplied research summary reports that distinct graphs can share strata and mentions \(B_1(K_3)\) versus \(B_1(G_{1,2-3})\), but the latter graph and that comparison are not defined here. It reports no located complete published characterization.
\subsection{Short English statement}
Characterize subgraphs of the complete graph that induce identical tensegrity strata on the intended common configuration space. The dimension and the stress-fiber/stratum equivalence must be fixed before graph equality criteria can be stated.
\subsection{Sources}
\begin{itemize}
\item Local UnsolvedMath registry, permanent ID \(6900004\), exact source code \texttt{catalogue entry}, ``Configuration Spaces of Tensegrities --- Problem 4.'' Proposer Oleg Karpenkov, 2020.
\item Karpenkov, \emph{Open Problems on Configuration Spaces of Tensegrities} (2020), Problem 4, recorded PDF page 4; \url{https://livrepository.liverpool.ac.uk/3022838/1/stratum.pdf}. Recorded access 2026-07-29; pdf-text-normalized extraction.
\item Supplied OPEN-TRIAGE summary refers to the source's \(B_1(K_3)\) and \(B_1(G_{1,2-3})\) discussion but provides no definition or separate characterization reference.
\end{itemize}
\end{document}
